\section{Terminology}

\textit{Scaling predictability}: "evaluation results from small-scale models reliably predict ablation effects in their larger-scale counterparts".

\begin{table*}[t]
    \centering
    \begin{tabular}{lll}
        \toprule
        \textbf{Symbol} & \textbf{Meaning} & \textbf{Options} \\
        \hline
        \(N\) & Non-embedding parameter count & 90M, 175M, 350M, 600M, 1B, 1.7B, 3B \\
        \hline
        \(A\) & Architecture & deep or shallow \\
        \hline
        \(O\) & Optimizer & AdEMAMix or Muon \\
        \hline
        \(A_f\) & Activation function & xielu or swiglu \\
        \hline
        \(K\) & Number of training languages & 1, 2, 8, 15, 30, 50 \\
        \hline
        \(M\) & Language-selection scheme & scheme A, B, or C \\
        \hline
        \(T\) & Sampling temperature & 1 or 3 \\
        \hline
        \(S\) & Random seed & 1904 (default), 64 and 313 (175M, 600M), 28 and 1797 (1B) \\
        \hline
        \(f\) & Fraction of training completed & 0.10, 0.20, ..., 0.70, 0.80, 0.85, 0.90, 0.95, 1.00 \\
        \bottomrule
        \toprule
        \textbf{Symbol} & \textbf{Meaning} & \textbf{Example} \\
        \hline
        \(b\) & Benchmark & HellaSwag \\
        \hline
        \(\ell\) & Evaluation language & Spanish \\
        \hline
        \(b_\ell\) & Language task & HellaSwag-es \\
        \hline
        \({b_\ell}_i\) & Sample/example & One instance inside HellaSwag-es \\
        \hline
        \(Y_{b_\ell}\) & Raw task score & accuracy or BPB \\
        \hline
        \(B\) & Evaluation suite & Complete collection of benchmarks \\
        \bottomrule
    \end{tabular}
    \caption{Notation symbols}
    \label{tab:notation-symbols}
\end{table*}

